\DocumentMetadata{
  pdfversion=2.0,
  pdfstandard=ua-2,
  lang=en-US,
  tagging=on,
  tagging-setup={math/setup=mathml-SE,tabsorder=structure}
}
\documentclass[11pt,letterpaper]{article}
\usepackage[margin=0.68in,includefoot]{geometry}
\usepackage{newcomputermodern}
\usepackage{amsmath,amscd,mathtools}
\usepackage{tikz,adjustbox}
\providecommand{\square}{\mdlgwhtsquare}
\usepackage[hidelinks]{hyperref}
\hypersetup{
  pdftitle={Algebra First-Year Exam, May 2015, Part 1},
  pdfauthor={University of Florida Department of Mathematics},
  pdfsubject={Graduate mathematics examination},
  pdfdisplaydoctitle=true,
  bookmarksopen=true
}
\setcounter{secnumdepth}{0}
\setlength{\parindent}{0pt}
\setlength{\parskip}{0.28em}
\setlength{\emergencystretch}{3em}
\setlength{\leftmargini}{2.15em}
\setlength{\leftmarginii}{2.65em}
\setlength{\leftmarginiii}{3em}
\renewcommand{\labelenumi}{\textbf{\theenumi.}}
\pagestyle{empty}
\raggedbottom
\allowdisplaybreaks
\newenvironment{examproblems}[1]{%
  \begin{enumerate}%
  \setcounter{enumi}{#1}%
  \setlength{\topsep}{0.35em}%
  \setlength{\partopsep}{0pt}%
  \setlength{\itemsep}{0.48em}%
  \setlength{\parsep}{0.18em}%
}{\end{enumerate}}
\newenvironment{examsubparts}{%
  \begin{enumerate}%
  \setlength{\topsep}{0.18em}%
  \setlength{\partopsep}{0pt}%
  \setlength{\itemsep}{0.16em}%
  \setlength{\parsep}{0.1em}%
}{\end{enumerate}}
\newenvironment{examinstructions}{%
  \par\begingroup\itshape
}{%
  \par\endgroup\vspace{0.18em}
}
\makeatletter
\renewcommand\section{\@startsection{section}{1}{0pt}%
  {-1.25ex plus -.25ex minus -.1ex}%
  {0.55ex plus .1ex}%
  {\normalfont\large\bfseries}}
\renewcommand{\@maketitle}{%
  \newpage\null\vskip -1em%
  {\footnotesize\bfseries
    \href{https://www.ufl.edu/}{University of Florida}, \href{https://math.ufl.edu/}{Department of Mathematics}\par}%
  \vskip -0.5em%
  \tagstructbegin{tag=Title}%
    {\LARGE\bfseries\href{https://gma.math.ufl.edu/exams/algebra-first-year/}{Algebra First-Year Exam}, May 2015, Part 1\par}%
  \tagstructend%
  \vskip 0.25em%
  \tagmcbegin{artifact}%
    \hrule height 1pt%
  \tagmcend%
  \vskip 1em%
}
\makeatother
\title{Algebra First-Year Exam, May 2015, Part 1}
\author{}
\date{}
\begin{document}
\maketitle
\thispagestyle{empty}
\begin{examinstructions}
Answer four problems. If you turn in more than four, only the first four will be graded. Within reason, you may use theorems as long as you state them clearly.
\end{examinstructions}
\begin{examproblems}{0}
\item
Let \(P \subset S_7\) be a cyclic subgroup of order 7. Show that the normalizer~\(N\) of~\(P\) has order 42, and find a pair of cycles generating~\(N\).
\item
Let \(Q_8\) be the quaternion group of order 8. Show that if \(Q_8\) acts faithfully on a set~\(S\), then \(|S| \geq 8\). Hint: the center lies in every nontrivial subgroup.
\item
Let~\(G\) be a group of order \(385=5\cdot 7\cdot 11\).
\begin{examsubparts}
\item[\textnormal{(i)}]
Show that the 7-Sylow and 11-Sylow subgroups are normal.
\item[\textnormal{(ii)}]
Show that the 7-Sylow subgroup is in the center.
\item[\textnormal{(iii)}]
Show that~\(G\) has a cyclic normal subgroup of order 77.
\end{examsubparts}
\item
This problem has two parts.
\begin{examsubparts}
\item[\textnormal{(i)}]
Find representatives for all conjugacy classes of elements of order 15 in \(S_{11}\).
\item[\textnormal{(ii)}]
Do the same for \(A_{11}\).
\end{examsubparts}
\item
Find
\begin{examsubparts}
\item[\textnormal{(i)}]
representatives of the isomorphism classes of abelian groups of order \(504=2^3 3^2 7\), and
\item[\textnormal{(ii)}]
the number of elements of order 6 in each such group.
\end{examsubparts}
\item
Suppose~\(G\) is a finite group, \(P\) is a~\(p\)-Sylow subgroup of~\(G\) for some prime~\(p\), and~\(N\) is a normal subgroup of~\(G\). Show that \(N\cap P\) is a~\(p\)-Sylow subgroup of~\(N\).
\end{examproblems}
\end{document}
